\begin{helper}
Let \(s,k\in\mathbb N\) with \(4\le s\le k\). Then
\[
  \sum_{t=1}^{s-1}\binom{k}{t}
    2^{(s-1)(t+k)-\binom{t+1}{2}}
  \le
  2^{sk+s^2/2},
\]
where the exponent on the left is interpreted as the natural-number exponent
displayed there, and the powers are viewed as real numbers.
\end{helper}
\begin{proof}
Fix \(t\in\{1,\ldots,s-1\}\). The binomial coefficient satisfies
\[
  \binom{t+1}{2}=\frac{t(t+1)}2.
\]
Since \(t\le s-1\), we have \(t+1\le 2(s-1)\), and hence
\[
  \binom{t+1}{2}\le (s-1)(t+k).
\]
Thus the natural subtraction in the exponent is ordinary subtraction after
casting to \(\mathbb R\).

Also,
\[
  \binom{t+1}{2}=\frac{t(t+1)}2\ge \frac{t^2}{2}.
\]
Therefore
\[
\begin{aligned}
  (s-1)(t+k)-\binom{t+1}{2}
  &\le (s-1)k+(s-1)t-\frac{t^2}{2} \\
  &\le (s-1)k+\frac{s^2}{2}.
\end{aligned}
\]
The last inequality follows from
\((s-t)^2\ge0\), which gives \(st-t^2/2\le s^2/2\), and from
\((s-1)t\le st\). Since \(2>1\), monotonicity of \(x\mapsto 2^x\) gives
\[
  2^{(s-1)(t+k)-\binom{t+1}{2}}
  \le 2^{(s-1)k+s^2/2}.
\]
Multiplying by the nonnegative \(\binom{k}{t}\) and summing over
\(1\le t\le s-1\), we get
\[
  \sum_{t=1}^{s-1}\binom{k}{t}
    2^{(s-1)(t+k)-\binom{t+1}{2}}
  \le
  2^{(s-1)k+s^2/2}
  \sum_{t=1}^{s-1}\binom{k}{t}.
\]
Because \(s\le k\), the index set \(\{1,\ldots,s-1\}\) is contained in
\(\{0,\ldots,k\}\). Hence Pascal's identity for the sum of a row gives
\[
  \sum_{t=1}^{s-1}\binom{k}{t}
  \le
  \sum_{t=0}^{k}\binom{k}{t}
  =
  2^k.
\]
Combining the previous two displays yields
\[
  \sum_{t=1}^{s-1}\binom{k}{t}
    2^{(s-1)(t+k)-\binom{t+1}{2}}
  \le
  2^k\cdot 2^{(s-1)k+s^2/2}
  =
  2^{sk+s^2/2},
\]
which is the desired estimate.
\end{proof}
